\subsection{Induced measure on a measurable subspace}

Let X be a subset of T, and let \(\nu\) be the restriction of the mapping \(\mu: K \mapsto \mu_{K}\) to the set of compact subsets of X; it is clear that \(\nu\) is a premeasure on X. On the other hand, let \(x \in X\) and let V be an open neighborhood of \(x\) in \(\mathbf{T}\) such that \(\mu^{\bullet}(\mathrm{V}) < +\infty\) ; then

\[
\nu^{\bullet}(\mathrm{X}\cap \mathrm{V}) = \sup_{\substack{\mathrm{K\ compact}\\ \mathrm{K}\subset \mathrm{X}\cap \mathrm{V}}}\mu^{\bullet}(\mathrm{K})\leqslant \mu^{\bullet}(\mathrm{V}) <   + \infty
\]

by Remark 3 of §1, No.~2, so that \(\nu\) is a measure.

When X is not \(\mu\) -measurable, the encumbrances \(\nu^{\bullet}\) and \((\mu^{\bullet})_{X}\) are not necessarily equal and the measure \(\nu\) presents no interest.

DEFINITION 1. --- Let X be a \(\mu\) -measurable subset of T. The restriction of \(\mu: \mathrm{K} \mapsto \mu_{\mathrm{K}}\) to the set of compact subsets of X is called the measure induced by \(\mu\) on the subspace X, and is denoted by \(\mu_{\mathrm{X}}\) or \(\mu|_{\mathrm{X}}\) .

PROPOSITION 1. --- Let X be a \(\mu\) -measurable subset of T. The encumbrance \((\mu_{\mathrm{X}})^{\bullet}\) is equal to the encumbrance \((\mu^{\bullet})_{\mathrm{X}}\) induced by \(\mu^{\bullet}\) on X (§1, No.~1). In other words, \((\mu_{\mathrm{X}})^{\bullet}(g) = \mu^{\bullet}(g^{0})\) for every function \(g \in \mathcal{F}_{+}(\mathrm{X})\) .

Let \(f \in \mathcal{F}_{+}(X)\) and let \(f^{0}\) be the extension by 0 of \(f\) to \(T\) . One has \((\mu^{\bullet})_{X}(f) = \mu^{\bullet}(f^{0}) = \sup_{L} \mu^{\bullet}(f^{0}\varphi_{L})\) , where \(L\) runs over the set of compact subsets of \(T\) (§1, No.~2, Prop. 2); similarly \((\mu_{X})^{\bullet}(f) = \sup_{K} \mu_{K}^{\bullet}(f_{K}) = \sup_{K} \mu^{\bullet}(f^{0}\varphi_{K})\) , where \(K\) runs over the set of compact subsets of \(X\) . Thus it all comes down to showing that \(\mu^{\bullet}(f^{0}\varphi_{L}) = \sup_{K} \mu^{\bullet}(f^{0}\varphi_{K})\) for every compact subset \(L\) of \(T\) , where \(K\) runs over the set of compact subsets of \(L \cap X\) . Now, let \((K_{n})\) be an increasing sequence of compact sets contained in \(L \cap X\) , such that \((L \cap X) - \bigcup_{n} K_{n}\) is locally \(\mu\) -negligible (§1, No.~8, Prop. 11); \(f^{0}\) being zero outside \(X\) , \(f^{0}\varphi_{L}\) is zero outside \(L \cap X\) hence is equal locally almost everywhere to the upper envelope of the sequence \((f^{0}\varphi_{K_{n}})\) . This implies that \(\mu^{\bullet}(f^{0}\varphi_{L}) = \sup \mu^{\bullet}(f^{0}\varphi_{K_{n}})\) , whence the desired result.

Remark 3). --- Let X be a \(\mu\) -measurable subset of T. By Prop. 10 of §1, No.~8, applied to \(g = \varphi_{\mathbf{X}}\) , there exists a crushing \((\mathrm{K}_{\alpha})_{\alpha \in \mathbf{A}}\) of T such that for each \(\alpha \in \mathbf{A}\) , either \(\mathrm{K}_{\alpha} \subset \mathrm{X}\) or \(\mathrm{K}_{\alpha} \subset \mathbb{C}\mathrm{X}\) . If one modifies the topology of T by the procedure of the Scholium of §1, No.~8, the space \(\mathrm{X}'\) obtained by equipping X with the topology induced by \(\mathrm{T}'\) is locally compact, and one knows how to associate to \(\mu\) (resp. to \(\mu_{\mathbf{X}}\) ) a measure \(\mu'\) (resp. \(\nu\) ) on \(\mathrm{T}'\) (resp. on \(\mathrm{X}'\) ) that admits the same essential upper integral as \(\mu\) (resp. \(\mu_{\mathbf{X}}\) ): this implies that \(\mu'_{\mathbf{X}'} = \nu\) . Since the locally negligible sets, measurable mappings, and essentially integrable functions with values in a Banach space, are the same for \(\mu\) and \(\mu'\) , and for \(\mu_{\mathbf{X}}\) and \((\mu_{\mathbf{X}})' = \nu = \mu'_{\mathbf{X}'}\) , the theory of integration with respect to an induced measure reduces to that which was treated in Ch. V, §7 in the special case of locally compact spaces. We leave to the reader the task of transcribing the results.
% semantic-fragments: legacy-input-seam
\relax{}
